\documentclass[troisieme]{fiche}
\metadonnees{16}{Le premier matin}{Bloc D — Rapporter et supposer}

\begin{document}
\entete

\objectifs{%
\textbf{Langue} : le passif au présent et au prétérit ; traduire \og on \fg.\par
\textbf{Texte} : E. Nesbit, \emph{The Railway Children} (1906), ch. \textsc{ii}.\par
\textbf{Durée} : 1 h — exercices 20 min, texte et questions 25 min, version 15 min.}

%% ===================================================================
\exercice[12 min]{Mettre au passif}

\rappel{\textbf{Rappel.} Le passif, c'est \emph{be} au temps voulu + participe passé.
L'agent, quand on le donne, vient après \emph{by}.}

\consigne{Mettez les phrases au passif.}

\begin{anglais}
\begin{enumerate}[label=\arabic*.,itemsep=3mm,leftmargin=*]
  \item Aunt Emma aired all the sheets.
    \lignes{1}
    \corr{All the sheets were aired by Aunt Emma.}
  \item The men moved the furniture.
    \lignes{1}
    \corr{The furniture was moved by the men.}
  \item Mother lights the candle every evening.
    \lignes{1}
    \corr{The candle is lit by Mother every evening.}
  \item Somebody has broken the kettle.
    \lignes{1}
    \corr{The kettle has been broken.}
  \item They mistook the door for a cupboard's.
    \lignes{1}
    \corr{The door was mistaken for a cupboard's.}
  \item A sudden noise frightened the children.
    \lignes{1}
    \corr{The children were frightened by a sudden noise.}
  \item Somebody stole the coal from the heap.
    \lignes{1}
    \corr{The coal was stolen from the heap.}
  \item They set the table before eight.
    \lignes{1}
    \corr{The table was set before eight.}
  \item The rain spoiled the garden.
    \lignes{1}
    \corr{The garden was spoiled by the rain.}
  \item People call the house Three Chimneys.
    \lignes{1}
    \corr{The house is called Three Chimneys.}
\end{enumerate}
\end{anglais}

\note[Ce qui change]{Le complément d'objet devient sujet, le verbe passe à \emph{be} +
participe passé au temps de départ, et le sujet de la phrase active devient l'agent. Le temps
ne change pas : présent (3, 10), prétérit (1, 2, 5, 6, 7, 8, 9), present perfect (4).}

\note[Quand l'agent disparaît]{Si le sujet actif est \emph{somebody}, \emph{people},
\emph{they} au sens indéfini, on ne le reprend pas : \emph{the kettle has been broken} (4),
\emph{the coal was stolen} (7), \emph{the house is called} (10). C'est même la raison
principale d'employer le passif.}

%% ===================================================================
\exercice[8 min]{Traduire \og on \fg}

\rappel{\textbf{Rappel.} Le \og on \fg{} français n'a pas d'équivalent anglais : on passe par
le passif, ou par \emph{people}, \emph{you}, \emph{they}.}

\consigne{Traduisez en anglais.}

\begin{anglais}
\begin{enumerate}[label=\arabic*.,itemsep=3mm,leftmargin=*]
  \item On a construit la maison en 1850.
    \lignes{1}
    \corr{The house was built in 1850.}
  \item Ici on parle anglais.
    \lignes{1}
    \corr{English is spoken here.}
  \item La lettre a été écrite hier.
    \lignes{1}
    \corr{The letter was written yesterday.}
  \item On a retrouvé la clef.
    \lignes{1}
    \corr{The key has been found.}
  \item Le dîner sera servi à huit heures.
    \lignes{1}
    \corr{Dinner will be served at eight.}
  \item Le bruit ne fut jamais expliqué.
    \lignes{1}
    \corr{The noise was never explained.}
\end{enumerate}
\end{anglais}

\note[Les participes irréguliers]{Le passif oblige à les savoir : \emph{build → built},
\emph{speak → spoken}, \emph{write → written}, \emph{find → found}. Une erreur de participe
se voit immédiatement.}

%% ===================================================================
\newpage
\section{Texte}

\noindent\small\textit{Le père des enfants a été emmené un soir par deux messieurs, et
personne ne dit pourquoi. La mère a quitté Londres avec Roberta, Peter et Phyllis pour une
maison isolée à la campagne, sans domestiques. Ils y sont arrivés de nuit.}\normalsize

\begin{texte}
Next morning Roberta woke Phyllis by pulling her hair gently, but quite enough for her
purpose.

``Wake up! wake up!'' said Roberta. ``We're in the new house --- don't you remember? No
servants or anything. Let's get up and begin to be useful. We'll just creep down
mouse-quietly, and have everything beautiful before Mother gets up. I've woke Peter. He'll be
dressed as soon as we are.''

So they dressed quietly and quickly. Of course, there was no water in their room, so when
they got down they washed as much as they thought was necessary under the
\gl[a spout]{spout}{un bec (de pompe)} of the pump in the \gl[a yard]{yard}{une cour}. One
\gl[to pump]{pumped}{pomper} and the other washed. It was \gl[splashy]{splashy}{plein
d'éclaboussures} but interesting.

``It's much more fun than \gl[a basin]{basin}{une cuvette} washing,'' said Roberta. ``How
\gl[sparkly]{sparkly}{scintillant} the \gl[a weed]{weeds}{une mauvaise herbe} are between the
stones, and the \gl[moss]{moss}{la mousse} on the roof --- oh, and the flowers!''

The roof of the back kitchen sloped down quite low. It was made of
\gl[thatch]{thatch}{le chaume} and it had moss on it, and house-leeks and
\gl[stonecrop]{stonecrop}{l'orpin} and \gl[a wallflower]{wallflowers}{une giroflée}, and even
a \gl[a clump]{clump}{une touffe} of purple flag-flowers, at the far corner.

``This is far, far, far and away prettier than Edgecombe Villa,'' said Phyllis. ``I wonder
what the garden's like.''

``We mustn't think of the garden yet,'' said Roberta, with earnest energy. ``Let's go in and
begin to work.''

They lighted the fire and put the \gl[a kettle]{kettle}{une bouilloire} on, and they arranged
the \gl[crockery]{crockery}{la vaisselle} for breakfast; they could not find all the right
things, but a glass \gl[an ash-tray]{ash-tray}{un cendrier} made an excellent
\gl[a salt-cellar]{salt-cellar}{une salière}, and a newish \gl[a baking-tin]{baking-tin}{un
moule à gâteau} seemed as if it would do to put bread on, if they had any.

When there seemed to be nothing more that they could do, they went out again into the fresh
bright morning. ``We'll go into the garden now,'' said Peter. But somehow they couldn't find
the garden. They went round the house and round the house.

So they all sat down on a great flat grey stone that had pushed itself up out of the grass;
and when Mother came out to look for them at eight o'clock, she found them deeply asleep in a
contented, sun-warmed \gl[a bunch]{bunch}{un tas, une grappe}.

They had made an excellent fire, and had set the kettle on it at about half-past five. So
that by eight the fire had been out for some time, the water had all boiled away, and the
bottom was burned out of the kettle. Also they had not thought of washing the crockery before
they set the table.
\end{texte}
\source{E. Nesbit, \emph{The Railway Children}, Londres, Wells Gardner, 1906,
ch.~\textsc{ii}.}

\partiel[15 min]{Questions}
\consigne{Répondez aux deux premières questions en français, aux deux dernières en anglais.}

\begin{enumerate}[label=\arabic*.,itemsep=2.5mm,leftmargin=*]
  \item Que décident de faire les enfants en se levant, et pourquoi ?
    \lignes{2}
    \corr{Descendre sans bruit et tout préparer avant que leur mère se lève : il n'y a plus de
    domestiques, et ils veulent se rendre utiles.}
  \item Qu'est-ce qui a mal tourné pendant qu'ils dormaient dehors ?
    \lignes{2}
    \corr{Le feu allumé à cinq heures et demie s'est éteint, l'eau de la bouilloire s'est
    entièrement évaporée et le fond de la bouilloire a brûlé. Ils avaient en outre mis la
    table sans laver la vaisselle.}
  \item \en{The washing under the pump is ``splashy but interesting''. What does that tell us
    about the way the children take their new life?}
    \lignes{3}
    \corr{They turn every discomfort into a discovery: no water in the room becomes a game,
    the mossy roof becomes a garden. The narrator never says they are poor now; she shows
    children who have not yet understood it.}
  \item \en{Find the two passive forms in the text and say why no agent is given.}
    \lignes{3}
    \corr{\emph{It was made of thatch} et \emph{the bottom was burned out of the kettle}.
    Personne n'a fait fondre le toit, et pour la bouilloire l'agent serait la maladresse des
    enfants : la narratrice ne la nomme pas, elle laisse le résultat parler. Attention,
    \emph{the water had all boiled away} n'est pas un passif : \emph{boil} y est intransitif.}
\end{enumerate}

%% ===================================================================
\section{Version}
\consigne{Traduisez le passage en français. Il suit le texte : la mère vient de réveiller les
enfants.}

\begin{version}
``But it doesn't matter --- the cups and saucers, I mean,'' said Mother. ``Because I've found
another room --- I'd quite forgotten there was one. And it's magic! And I've boiled the water
for tea in a \gl[a saucepan]{saucepan}{une casserole}.''

The forgotten room opened out of the kitchen. In the agitation and half darkness the night
before its door had been mistaken for a \gl[a cupboard]{cupboard}{un placard}'s. It was a
little square room, and on its table, all nicely set out, was a
\gl[a joint]{joint}{un rôti, une pièce de viande} of cold \gl[roast]{roast}{rôti} beef, with
bread, butter, cheese, and a pie.

``Pie for breakfast!'' cried Peter; ``how perfectly \gl[ripping]{ripping}{épatant,
formidable}!''

``It isn't pigeon-pie,'' said Mother; ``it's only apple. Well, this is the supper we ought to
have had last night.''

That was a wonderful breakfast. It is unusual to begin the day with cold apple pie, but the
children all said they would rather have it than meat.
\end{version}
\source{\emph{Ibid.}}

\lignes{16}

\corr{\textbf{Proposition de traduction.} \og Mais cela n'a pas d'importance --- les tasses et
les soucoupes, je veux dire, dit Maman. Parce que j'ai trouvé une autre pièce --- j'avais
complètement oublié qu'il y en avait une. Et elle est magique ! Et j'ai fait bouillir l'eau du
thé dans une casserole. \fg{} La pièce oubliée donnait sur la cuisine. Dans l'agitation et la
demi-obscurité de la veille au soir, on avait pris sa porte pour celle d'un placard. C'était
une petite pièce carrée, et sur sa table, joliment disposés, il y avait un rôti de bœuf
froid, du pain, du beurre, du fromage et une tourte. \og De la tourte au petit déjeuner !
s'écria Peter. Comme c'est épatant ! \fg{} \og Ce n'est pas de la tourte au pigeon, dit
Maman, seulement de la tourte aux pommes. Enfin, c'est le souper que nous aurions dû prendre
hier soir. \fg{} Ce fut un merveilleux petit déjeuner. Il est rare de commencer la journée par
de la tourte aux pommes froide, mais les enfants dirent tous qu'ils la préféraient à de la
viande.}

\note[Choix de traduction 1]{\emph{its door had been mistaken for a cupboard's} : passif sans
agent, et l'agent serait ici toute la famille. Le français passe par \og on \fg{} — \og on
avait pris sa porte pour \fg.}

\note[Choix de traduction 2]{\emph{all nicely set out, was a joint of cold roast beef} :
l'anglais rejette le sujet après le verbe pour faire attendre la découverte. Le français
rétablit l'ordre normal avec \og il y avait \fg.}

\note[Choix de traduction 3]{\emph{they would rather have it than meat} : \emph{would rather}
= \og préférer \fg, suivi de la base verbale sans \emph{to}. À ne pas confondre avec
\emph{had rather} ni avec \emph{had better}.}

%% ===================================================================
\partiel{Lexique à mémoriser}
\begin{multicols}{2}\small
\begin{itemize}[label=--,itemsep=0.8mm,leftmargin=*]
  \item \textbf{a sheet} — un drap ; \textit{et} une feuille
  \item \textbf{a kettle} — une bouilloire
  \item \textbf{a saucepan} — une casserole
  \item \textbf{crockery} — la vaisselle
  \item \textbf{a cupboard} — un placard
  \item \textbf{to boil} — bouillir, faire bouillir
  \item \textbf{to burn} — brûler ; burnt, burnt
  \item \textbf{a pump} — une pompe
  \item \textbf{moss} — la mousse (végétale)
  \item \textbf{quietly} — sans bruit
  \item \textbf{a mistake} — une erreur \textit{(fiche 9)}
  \item \textbf{breakfast} — le petit déjeuner
\end{itemize}
\end{multicols}

\end{document}
